%%%%%%%%%%%%%%%%%%%%%%%%%%%%%%%%%%%%%%%%%%%%%%%%%%%%%%%%%%%%%%%%%%%%%%%%%%%%%%%
%% Capítulo 1 --- Introdução
%%
%% Citações: abntex2cite (sistema autor-data). \cite{} → (AUTOR, ano);
%%           \citeonline{} → Autor (ano).
%%
%% MARCAÇÃO DE TRECHO EM ANDAMENTO --------------------------------------------
%% \emandamento{...} destaca em vermelho os trechos cuja validade muda quando o
%% protótipo for construído e ensaiado em ambiente real. É marcação de rascunho.
%%
%% Para REMOVER todas as marcas na versão final, troque a definição abaixo por:
%%     \newcommand{\emandamento}[1]{#1}
%% ou apague as chamadas \emandamento{...} do texto.
%%%%%%%%%%%%%%%%%%%%%%%%%%%%%%%%%%%%%%%%%%%%%%%%%%%%%%%%%%%%%%%%%%%%%%%%%%%%%%%

\makeatletter
\@ifundefined{emandamento}{%
  \@ifundefined{textcolor}%
    {\newcommand{\emandamento}[1]{#1}}%
    {\newcommand{\emandamento}[1]{\textcolor{red}{#1}}}%
}{}
\makeatother

\chapter{Introdução}
\label{cap:introducao}

A manutenção preditiva de máquinas rotativas por análise de vibração é uma prática
normalizada. A norma ISO 20816-3:2022 estabelece faixas de severidade a partir da
velocidade eficaz de vibração de banda larga e define os limites que separam operação
aceitável de operação inaceitável para máquinas industriais acima de
\SI{15}{\kilo\watt}, com rotação entre \num{120} e \num{30000}~r/min. Para essa medição,
a norma exige do instrumento resposta plana em faixa de \emph{ao menos}
\SIrange{10}{1000}{\hertz}~\cite{iso208163_2022}. O procedimento é
consolidado, e o obstáculo à sua adoção não é metodológico: é de custo de instrumentação.
Acelerômetros piezoelétricos com a banda exigida para diagnóstico, seus condicionadores e
os analisadores associados restringem a cobertura aos ativos considerados críticos, e
deixam sem instrumentação a maior parte do parque de motores de uma planta.

A resposta que a literatura recente oferece a esse obstáculo é o nó sensor de baixo custo:
uma unidade de medição inercial MEMS acoplada a um microcontrolador de baixo consumo, com
detecção de anomalia executada no próprio dispositivo e enlace de rádio de longo alcance
para telemetria. A viabilidade computacional dessa arquitetura está estabelecida. O
algoritmo Isolation Forest~\cite{liu2008} treina e executa dentro de um
ESP32~\cite{antonini2023}, existem variantes desenhadas para reduzir seu custo em
TinyML~\cite{barbariol2022}, e a combinação de microcontrolador, acelerômetro MEMS de
consumo e detecção não supervisionada já foi demonstrada em manutenção preditiva de baixo
custo~\cite{kolok2025}. O que a viabilidade computacional não resolve --- e é o ponto de
partida deste trabalho --- é a questão anterior: o que uma cadeia de medição dessa classe
efetivamente permite medir.

Essa questão tem uma ordem de precedência que organiza todo o restante. Primeiro vem a
aquisição: piso de ruído, banda passante e caminho mecânico entre a fonte de vibração e o
elemento sensor. Nenhuma decisão posterior recupera informação que o sensor não coletou.
Depois vem a estimação, que decide o que se faz com a informação disponível: a política de
limiar, a métrica de desempenho e o protocolo de validação. Por último vem a implantação
--- energia, comportamento térmico e invólucro ---, que determina por quanto tempo o
dispositivo continua operando. Os três elos são encadeados, e um erro no primeiro não é
compensável nos seguintes.

É contra essa cadeia que se identifica a lacuna. A literatura de vibração embarcada é
forte no primeiro elo e frágil no segundo. Sobre aquisição, há consenso quantitativo:
nenhum estudo de adequação de sensor revisado testou dispositivo com limite superior
abaixo de \SI{1500}{\hertz}~\cite{albarbar2008}, e os grupos que se propõem a diagnosticar
defeito de rolamento com MEMS escolhem sensores de dezenas de
quilohertz~\cite{landi2023,jakobsen2024}, porque a assinatura do defeito viaja modulada em
uma portadora de alta frequência~\cite{randall2011,antoni2007}. Sobre estimação, ao
contrário, o registro é lacunar. O escore do Isolation Forest tem \num{0,5} como
referência interpretativa, e não como corte calibrado~\cite{liu2008}; a implementação de
referência transporta essa interpretação para um deslocamento fixo, sem alterar o
ranqueamento~\cite{scikitlearn2026}; e os trabalhos que reportam bom desempenho com o
detector abandonam o padrão sem que a política adotada seja submetida a
varredura~\cite{antonini2023,brito2022}. Não se localizou, entre os trabalhos que executam
o detector em microcontrolador, nenhum que variasse o parâmetro de contaminação sobre a
mesma base e reportasse a fração de operação normal marcada como anomalia, nem taxa de
falso alarme publicada para detector treinado no próprio dispositivo --- que é justamente
o modo de operação exigido de uma máquina sem histórico rotulado. Soma-se a isso uma
fragilidade de protocolo com efeito documentado: janelas consecutivas de um mesmo ensaio
compartilham a assinatura da montagem, e sua divisão aleatória entre treino e teste
permite que o modelo identifique o ensaio em vez da falha. \citeonline{varejao2025}
denominam o fenômeno \emph{similarity bias} e relatam que, de \num{41} estudos revisados
sobre uma base de referência consolidada, \num{40} adotaram delineamento suscetível a ele.

Este trabalho ocupa essa lacuna com uma tese delimitada. Uma cadeia de medição de baixo
custo, com banda e piso de ruído declarados, sustenta a detecção de desbalanceamento sob taxa de falso
alarme fixada, mas não sustenta
diagnóstico de defeito de rolamento; e a contribuição verificável de um dispositivo dessa
classe está no segundo elo da cadeia --- limiar calibrado por risco alvo, métrica avaliada
no ponto de operação e protocolo de validação sem vazamento entre ensaios ---, contribuição
que só é válida se o primeiro elo estiver declarado em vez de disfarçado. O objeto de
estudo é o Kaelix, um dispositivo que combina acelerômetro MPU6050, microcontrolador
ESP32-S3, enlace LoRa e invólucro em ASA impresso por deposição, com detecção de anomalia
embarcada e rotulagem de severidade pela norma.

A validade dessa tese depende de que se declare a procedência de cada resultado, e ela não
é uniforme. Os resultados de detecção por modo de falha provêm de dado real: sinais de uma
bancada instrumentada, adquiridos a \SI{50}{\kilo\hertz} e submetidos a decimação
controlada em \emph{software} para emular a banda do dispositivo. O que é modelado, nesse
caso, não é o sinal, é a aquisição. Os resultados sobre banda, integração, limiar, métrica e
protocolo de validação provêm de sinal sintetizado, cuja resposta correta é conhecida, e
sustentam relações entre métodos, não desempenho em máquina. Os resultados de análise
mecânica e térmica provêm de modelagem computacional sobre as especificações de projeto,
tomadas como dadas. \emandamento{E a cadeia física do próprio
dispositivo não foi verificada: nenhum ensaio físico foi conduzido e nenhuma medição de
campo alimenta os modelos, uma vez que a construção do protótipo e seu ensaio em ambiente
real constituem a etapa seguinte deste trabalho.} Essa distinção restringe uma afirmação
específica e é declarada aqui por isso: \emandamento{os resultados sustentam conclusões
sobre o que a cadeia de aquisição especificada permite medir, e não sobre o desempenho do
dispositivo montado em operação industrial.} As duas coisas serão mantidas separadas ao
longo do texto, e cada resultado indicará de qual das três origens provém.

\section{Objetivo}
\label{sec:objetivo}

\subsection{Objetivo geral}
\label{ssec:objetivo-geral}

Determinar o que a cadeia de medição de vibração especificada para o Kaelix permite medir,
e calibrar o estágio de decisão do detector embarcado de modo que o desempenho reportado
seja verificável no ponto de operação pretendido.

\subsection{Objetivos específicos}
\label{ssec:objetivo-especifico}

Dentro desse objetivo geral, encontram-se como metas específicas:

\begin{itemize}
    \item Quantificar a perda de contraste do indicador de defeito de rolamento entre a
          cadeia de referência e a cadeia truncada do dispositivo, e verificar se o
          delineamento permite atribuir essa perda à banda ou ao método de demodulação;

    \item Determinar o método de integração de aceleração para velocidade compatível com a
          norma, e quantificar o erro dos métodos alternativos contra um sinal cuja
          integral é conhecida analiticamente~\cite{brandt2014,iso208163_2022};

    \item Medir a taxa de falso alarme produzida pela política de limiar padrão do
          Isolation Forest e o efeito do parâmetro de contaminação sobre ela;

    \item Comparar o desempenho do detector sob divisão aleatória por janela e sob divisão
          por ensaio, e quantificar a inflação de acurácia atribuível ao vazamento de
          assinatura de montagem;

    \item Avaliar, por modelagem, a resposta modal e a margem térmica do conjunto
          invólucro-sensor-fixação, e identificar qual componente impõe o limite de
          operação;

    \item Declarar, para cada questão levantada, se ela foi respondida por medição, por
          modelagem ou se permanece dependente de ensaio físico.
\end{itemize}

\section{Metodologia}
\label{sec:metodologia}

A estratégia adotada é comum às questões investigadas: construir uma condição de
referência cujo resultado correto decorre de solução analítica ou de definição normativa,
submetê-la à cadeia de medição proposta e medir o desvio introduzido pela configuração do
projeto. Quando a resposta certa é conhecida de antemão, um erro de implementação não
consegue se esconder atrás de um resultado plausível.

O trabalho se desenvolve em quatro etapas. A primeira é a revisão da literatura,
organizada por eixos correspondentes aos elos da cadeia descrita, com verificação de cada
referência em fonte primária e registro explícito das lacunas encontradas, incluindo
aquelas que a própria verificação derrubou. A segunda é a caracterização da aquisição: um sinal
sintetizado com resposta conhecida é submetido à cadeia de referência e à cadeia truncada, e
os descritores extraídos são comparados entre elas; em paralelo, sinais de bancada
adquiridos em alta taxa são decimados de forma controlada para emular a banda do
dispositivo e medir a detecção por modo de falha.
A terceira é a calibração do estágio de decisão, na qual a política de limiar é avaliada
por varredura, o desempenho é reportado no ponto de operação de baixa taxa de falso alarme
e o protocolo de validação é submetido à divisão por ensaio. A quarta é a verificação de
implantação, por modelagem mecânica e térmica do conjunto montado a partir das
especificações de projeto.

Duas decisões metodológicas atravessam as quatro etapas. A primeira é a rastreabilidade:
cada resultado é acompanhado da indicação de sua origem --- sinal sintetizado, dado real
decimado ou modelagem sobre especificação --- e os artefatos que o produzem, dados
intermediários e código de geração, são mantidos versionados e regeneráveis. A exceção são
os resultados de treino sobre dado real, que não têm artefato versionado próprio e exigem
nova execução do treino para serem reproduzidos (Seção~\ref{sec:dev-sinal}). A
segunda é o registro do que não foi verificado. As questões levantadas que dependem de
ensaio físico ou de consulta regulatória são enunciadas, mantidas abertas e associadas à
afirmação cujo alcance elas restringem, em vez de omitidas. Esse registro é parte do
método, e não uma ressalva sobre ele.

\section{Estrutura do Trabalho}
\label{sec:estrutura}

Este trabalho está organizado em seis capítulos. O Capítulo~\ref{cap:introducao}
apresentou o contexto, a lacuna identificada, a tese defendida e a procedência dos
resultados que a sustentam. O Capítulo~\ref{cap:fundamentacao} apresenta os conceitos
em que o trabalho se apoia: vibração de máquinas rotativas e suas assinaturas de falha,
aquisição e processamento de sinais, avaliação normativa de severidade, detecção de
anomalias e sua avaliação, e as restrições de um dispositivo embarcado. O
Capítulo~\ref{cap:revisao} revisa a literatura por eixos, seguindo a ordem
de precedência da cadeia de medição, e consolida o que a literatura sustenta, o que ela
contesta e quais lacunas permanecem abertas. O Capítulo~\ref{cap:desenvolvimento} descreve o desenvolvimento do
dispositivo: a cadeia de aquisição, o detector embarcado, o projeto de hardware e a
arquitetura de comunicação e recepção. O Capítulo~\ref{cap:resultados} apresenta os resultados e os confronta
com a literatura revisada, organizado pela ordem de precedência da cadeia: primeiro o que
a aquisição permite medir, depois o que a estimação decide com a informação disponível e,
por fim, o que a implantação impõe como limite de operação. O Capítulo~\ref{cap:conclusoes}
reúne as conclusões, delimita seu alcance e enuncia o trabalho futuro decorrente das
questões mantidas em aberto.
